The current results motivate two interpretations. First, learned camera geometry should not be judged only as metric calibration. A pose predictor optimized jointly with action can learn a representation of camera geometry that is useful for the downstream policy, even when it is not a drop-in replacement for calibrated extrinsics in every setting. This can explain why partial pose supervision can outperform a calibrated baseline in the current LIBERO-MV Spatial result.

Second, novel-view synthesis provides a training signal that is aligned with camera-shift robustness. The decoder cannot reconstruct a held-out view from a target camera query unless the observation encoder preserves information about object state, occlusion, and cross-view geometry. This does not require the latent representation to be an explicit point cloud or voxel map; it only requires the representation to be renderable enough to support the auxiliary task. The latest MimicGen results make this the stronger empirical story: the NVS variant improves tight-camera success over the action-only variant on every completed task.

The current draft should still remain conservative about explicit 3D comparisons. ManiFlow remains strong on some tasks, especially Stack tight-camera evaluation. The sharper claim is not that learned NVS representations uniformly dominate point clouds, but that end-to-end camera learning can use episodes without extrinsic labels: calibrated ray and point-cloud methods need known geometry to build their explicit features, whereas our policy can still receive action and NVS gradients through predicted geometry.
